\capitulo{Covariancia y correlaci\'on}{Medir si dos variables se mueven juntas, en sentidos opuestos o cada una por su lado.}{cap:cov}
\meta{85}{Ej. 3.12 \textbullet{} Serie 9 y 10 \textbullet{} Autoexamen 7 y 10}

\subsection{Qu\'e miden}

La variancia mide cu\'anto se desparrama \textbf{una} variable. La \textbf{covariancia} mide c\'omo se mueven \textbf{dos} variables
\textbf{una respecto de la otra}.

\begin{center}
\begin{tikzpicture}[scale=1.0]
\foreach \tit/\sub/\pts [count=\k from 0] in {
  {$\rho>0$}/{cuando $X$ sube, $Y$ tiende a subir}/{(0.3,0.5),(0.6,0.4),(0.8,0.9),(1.1,0.8),(1.3,1.3),(1.6,1.2),(1.8,1.7),(2.1,1.6),(2.3,2.1),(2.6,2.0),(1.0,1.2),(1.9,1.4)},
  {$\rho<0$}/{cuando $X$ sube, $Y$ tiende a bajar}/{(0.3,2.1),(0.6,2.2),(0.8,1.7),(1.1,1.9),(1.3,1.4),(1.6,1.5),(1.8,1.0),(2.1,1.1),(2.3,0.6),(2.6,0.7),(1.0,1.5),(1.9,1.2)},
  {$\rho\approx0$}/{sin relaci\'on lineal}/{(0.4,0.6),(0.5,1.9),(0.9,1.2),(1.0,2.2),(1.3,0.5),(1.5,1.6),(1.7,0.9),(2.0,2.1),(2.2,1.3),(2.5,0.6),(2.6,1.8),(1.2,1.5)}}{
  \begin{scope}[xshift=\k*5.4cm]
    \draw[eje] (0,0)--(3.0,0) node[right,et]{$x$}; \draw[eje] (0,0)--(0,2.7) node[above,et]{$y$};
    \foreach \p in \pts \fill[colY] \p circle (1.6pt);
    \node[font=\sffamily\bfseries\small,text=tinta] at (1.5,-0.45) {\tit};
    \node[font=\footnotesize,text=gris] at (1.5,-0.9) {\sub};
  \end{scope}}
\end{tikzpicture}
\end{center}

\begin{idea}{Covariancia, correlaci\'on y determinaci\'on}
\[
\Cov(X,Y)=\sigma_{XY}=\E(XY)-\E(X)\,\E(Y)
\qquad
\rho=\frac{\Cov(X,Y)}{\sigma_X\,\sigma_Y}
\qquad
\rho^2=\frac{[\Cov(X,Y)]^2}{\Var(X)\,\Var(Y)}
\]
\begin{itemize}
\item \textbf{Covariancia:} su \textbf{signo} dice el sentido de la relaci\'on. Su tama\~no depende de las unidades, por eso no se interpreta sola.
\item \textbf{Coeficiente de correlaci\'on $\rho$:} es la covariancia sin unidades. Siempre cumple $-1\le\rho\le1$ y mide qu\'e tan
\textbf{lineal} es la relaci\'on. Vale $\pm1$ s\'olo si $Y=mX+b$ exactamente.
\item \textbf{Coeficiente de determinaci\'on $\rho^2$:} fracci\'on de la variaci\'on de $Y$ que se explica conociendo $X$.
\item \textbf{Si son independientes:} $\Cov(X,Y)=0$ y $\rho=0$. Al rev\'es no vale: covariancia cero no garantiza independencia.
\end{itemize}
\end{idea}

\porque{La f\'ormula $\E(XY)-\E(X)\E(Y)$ compara el promedio del producto con el producto de los promedios.
Si fueran independientes ser\'ian iguales y la resta dar\'ia cero; lo que sobra es lo que tienen de relaci\'on.}

\begin{receta}{covariancia y correlaci\'on: las cinco esperanzas}
\begin{pasosreceta}
\item Saca las marginales y con ellas calcula $\E(X)$, $\E(X^2)$, $\E(Y)$, $\E(Y^2)$.
\item Calcula $\E(XY)$ con la conjunta (suma de $x\cdot y\cdot$celda, o integral doble de $xy\,\fxy$).
\item Variancias: $\Var(X)=\E(X^2)-[\E(X)]^2$ \ y \ $\Var(Y)=\E(Y^2)-[\E(Y)]^2$.
\item Covariancia: $\Cov(X,Y)=\E(XY)-\E(X)\,\E(Y)$.
\item Correlaci\'on: $\rho=\dfrac{\Cov(X,Y)}{\sqrt{\Var(X)}\sqrt{\Var(Y)}}$. Comprueba que quede entre $-1$ y $1$.
\end{pasosreceta}
\textbf{Atajo:} si ya sabes que son independientes, $\Cov=0$ y $\rho=0$ sin calcular nada.
\end{receta}

\begin{resuelto}{Una tabla peque\~na, completa}{Apuntes \textbullet{} Autoexamen 10}
\begin{enunciado}
\begin{center}\small
\begin{tabular}{|C|C|C|C|}\hline
\hd{$y\,\backslash\,x$} & \hx{0} & \hx{1} & \hy{$f_Y(y)$}\\\hline
\hy{1} & 0.1 & 0.3 & \my{0.4}\\\hline
\hy{2} & 0.4 & 0.2 & \my{0.6}\\\hline
\hx{$f_X(x)$} & \mx{0.5} & \mx{0.5} & \textbf{1}\\\hline
\end{tabular}
\end{center}
Obtener el coeficiente de correlaci\'on entre $X$ y $Y$.
\end{enunciado}
\begin{pasos}
\paso{Esperanzas de $X$ \textnormal{(con $f_X$)}}
\[\E(X)=0(0.5)+1(0.5)=0.5\qquad\E(X^2)=0^2(0.5)+1^2(0.5)=0.5\]
\paso{Esperanzas de $Y$ \textnormal{(con $f_Y$)}}
\[\E(Y)=1(0.4)+2(0.6)=1.6\qquad\E(Y^2)=1^2(0.4)+2^2(0.6)=0.4+2.4=2.8\]
\paso{$\E(XY)$ \textnormal{(con la conjunta; las celdas con $x=0$ aportan cero)}}
\[\E(XY)=(1)(1)(0.3)+(1)(2)(0.2)=0.3+0.4=0.7\]
\paso{Variancias}
\[\Var(X)=0.5-(0.5)^2=0.25\qquad\Var(Y)=2.8-(1.6)^2=2.8-2.56=0.24\]
\paso{Covariancia}
\[\Cov(X,Y)=0.7-(0.5)(1.6)=0.7-0.8=-0.1\]
\paso{Correlaci\'on}
\[\rho=\frac{-0.1}{\sqrt{0.25}\,\sqrt{0.24}}=\frac{-0.1}{(0.5)(0.4899)}=\res{-0.408}\]
\end{pasos}
\enpalabras{relaci\'on negativa de fuerza moderada: cuando $X$ vale 1, $Y$ tiende a valer 1 en lugar de 2. Es el inciso A del autoexamen.}
\end{resuelto}

\begin{resuelto}{Transistores}{Serie 3 \textbullet{} problema 10 b}
\begin{enunciado}
Con la tabla construida en el cap\'itulo \ref{cap:discretas}, calcular el coeficiente de correlaci\'on.
\begin{center}\small
\begin{tabular}{|C|C|C|C|C|}\hline
\hd{$x\,\backslash\,y$} & \hy{1} & \hy{2} & \hy{3} & \hx{$f_X(x)$}\\\hline
\hx{1} & 0.3 & 0.2 & 0.1 & \mx{0.6}\\\hline
\hx{2} & 0.2 & 0.1 & 0 & \mx{0.3}\\\hline
\hx{3} & 0.1 & 0 & 0 & \mx{0.1}\\\hline
\hy{$f_Y(y)$} & \my{0.6} & \my{0.3} & \my{0.1} & \textbf{1}\\\hline
\end{tabular}
\end{center}
\end{enunciado}
\begin{pasos}
\paso{Esperanzas de $X$}
\[\E(X)=1(0.6)+2(0.3)+3(0.1)=1.5\qquad\E(X^2)=1(0.6)+4(0.3)+9(0.1)=2.7\]
\paso{Esperanzas de $Y$}
Las marginales son iguales, as\'i que $\E(Y)=1.5$ y $\E(Y^2)=2.7$.
\paso{$\E(XY)$: cada celda por su $x$ y su $y$}
\begin{align*}
\E(XY)&=(1)(1)(0.3)+(1)(2)(0.2)+(1)(3)(0.1)+(2)(1)(0.2)+(2)(2)(0.1)+(3)(1)(0.1)\\
&=0.3+0.4+0.3+0.4+0.4+0.3=2.1
\end{align*}
\paso{Variancias}
\[\Var(X)=\Var(Y)=2.7-(1.5)^2=2.7-2.25=0.45\]
\paso{Covariancia}
\[\Cov(X,Y)=2.1-(1.5)(1.5)=2.1-2.25=-0.15\]
\paso{Correlaci\'on}
\[\rho=\frac{-0.15}{\sqrt{0.45}\,\sqrt{0.45}}=\frac{-0.15}{0.45}=\res{-\tfrac13}\]
\end{pasos}
\enpalabras{es negativa porque si el primer defectuoso tarda en aparecer, quedan pocas piezas y el segundo aparece pronto.}
\end{resuelto}

\begin{resuelto}{Neum\'aticos: caso continuo}{Serie 3 \textbullet{} problema 9}
\begin{enunciado}
$\fxy(x,y)=K(x^2+y^2)$ para $20\le x\le30$, $20\le y\le30$, con $K=\tfrac{3}{380000}$.\quad
a) Covariancia.\quad b) Coeficiente de correlaci\'on.
\end{enunciado}
Del cap\'itulo \ref{cap:margcond}: $f_X(x)=\dfrac{3x^2}{38000}+\dfrac1{20}$, y por simetr\'ia $f_Y$ tiene la misma forma. Todo lo de $Y$ vale lo mismo que lo de $X$.
\begin{pasos}
\paso{$\E(X)$ con la marginal}
\[\E(X)=\int_{20}^{30}x\Big(\tfrac{3x^2}{38000}+\tfrac1{20}\Big)dx=\Big[\tfrac{3x^4}{152000}+\tfrac{x^2}{40}\Big]_{20}^{30}
=\tfrac{3(810000-160000)}{152000}+\tfrac{900-400}{40}\]
\[=12.8289+12.5=25.3289\qquad\Big(\text{exacto: }\tfrac{1925}{76}\Big)\]
\paso{$\E(X^2)$ con la marginal}
\[\E(X^2)=\int_{20}^{30}x^2\Big(\tfrac{3x^2}{38000}+\tfrac1{20}\Big)dx=\Big[\tfrac{3x^5}{190000}+\tfrac{x^3}{60}\Big]_{20}^{30}
=\tfrac{3(24300000-3200000)}{190000}+\tfrac{27000-8000}{60}\]
\[=333.1579+316.6667=649.8246\]
\paso{$\E(XY)$ con la conjunta}
\[\E(XY)=K\int_{20}^{30}\!\!\int_{20}^{30}xy\,(x^2+y^2)\,dx\,dy=K\int_{20}^{30}\!\!\int_{20}^{30}\big(x^3y+xy^3\big)\,dx\,dy\]
Los dos t\'erminos valen lo mismo por simetr\'ia, y cada uno se separa en dos integrales simples:
\[\int_{20}^{30}x^3dx=\tfrac{810000-160000}{4}=162500\qquad\int_{20}^{30}y\,dy=\tfrac{900-400}{2}=250\]
\[\E(XY)=K\cdot2\,(162500)(250)=\tfrac{3}{380000}(81250000)=641.4474\]
\paso{Variancia}
\[\Var(X)=\Var(Y)=649.8246-(25.3289)^2=649.8246-641.5556=8.2690\]
\paso{Covariancia}
\[\Cov(X,Y)=641.4474-(25.3289)(25.3289)=641.4474-641.5556=\res{-0.1082}\]
\paso{Correlaci\'on}
\[\rho=\frac{-0.1082}{\sqrt{8.2690}\,\sqrt{8.2690}}=\frac{-0.1082}{8.2690}=\res{-0.0131}\]
\end{pasos}
\verifica{la serie da $-0.108$: es la covariancia. La correlaci\'on, $-0.0131$, indica una relaci\'on lineal pr\'acticamente nula.}
\begin{ojo}
\mal{Redondear $\E(X)$ a $25.33$: entonces $25.33^2=641.61$ y la covariancia sale $-0.16$, que est\'a mal.}
\bien{Aqu\'i se restan dos n\'umeros casi iguales. Trabaja con fracciones o con al menos cuatro decimales hasta el final.}
\end{ojo}
\end{resuelto}

\begin{resuelto}{Cuando son independientes no hay nada que calcular}{Apuntes \textbullet{} Ej.~3.12 c, d (en blanco en tus apuntes)}
\begin{enunciado}
$\fxy(x,y)=4xy\,e^{-(x^2+y^2)}$ para $x\ge0$, $y\ge0$.\quad c) Esperanzas condicionales.\quad d) Coeficiente de correlaci\'on.
\end{enunciado}
\begin{pasos}
\paso{Recupero lo que ya se sabe}
En el cap\'itulo \ref{cap:margcond} se prob\'o que son independientes, con $f_X(x)=2x\,e^{-x^2}$ y $f_Y(y)=2y\,e^{-y^2}$.
\paso{Inciso c: esperanzas condicionales}
Al ser independientes, las condicionales son las marginales, y las esperanzas condicionales son constantes:
\[\E[X\mid y]=\E(X)=\int_0^\infty x\cdot2x\,e^{-x^2}dx=\res{\tfrac{\sqrt\pi}{2}}\qquad\qquad\E[Y\mid x]=\E(Y)=\res{\tfrac{\sqrt\pi}{2}}\]
\paso{Inciso d: correlaci\'on}
Independientes $\Rightarrow\ \Cov(X,Y)=0\ \Rightarrow\ \res{\rho=0}$
\end{pasos}
\porque{La integral $\int_0^\infty2x^2e^{-x^2}dx=\tfrac{\sqrt\pi}2\approx0.886$ sale por partes usando $\int_0^\infty e^{-x^2}dx=\tfrac{\sqrt\pi}2$.
En el examen lo importante es el argumento: independencia implica esperanza condicional constante y correlaci\'on cero.}
\end{resuelto}

\begin{turno}
\textbf{1.} (Autoexamen 7) $\fxy=4xy$ para $0<x<1$, $0<y<1$. Halla $\Cov(X,Y)$.

\textbf{2.} Con la densidad de la tabacalera, $\fxy=24xy$ en $x+y\le1$ (marginales $f_X=12x(1-x)^2$ y $f_Y=12y(1-y)^2$), halla $\Cov(X,Y)$ y $\rho$.

\textbf{3.} Con $\fxy=x+y$ en el cuadrado unitario (Ej.~3.4), halla $\Cov(X,Y)$ y $\rho$.
\comprueba{\textbf{1.} soporte rectangular y $4xy=(2x)(2y)$: independientes, $\Cov=0$ (inciso C).\quad
\textbf{2.} $\E(X)=\int_0^112x^2(1-x)^2dx=12\big(\tfrac13-\tfrac12+\tfrac15\big)=\tfrac25$;
$\E(X^2)=12\big(\tfrac14-\tfrac25+\tfrac16\big)=\tfrac15$; $\Var(X)=\tfrac15-\tfrac4{25}=\tfrac1{25}$; lo mismo para $Y$.
$\E(XY)=\int_0^1\!\int_0^{1-x}24x^2y^2\,dy\,dx=\int_0^18x^2(1-x)^3dx=\tfrac2{15}$.
$\Cov=\tfrac2{15}-\tfrac4{25}=-\tfrac2{75}$; $\rho=\tfrac{-2/75}{1/25}=-\tfrac23$.\quad
\textbf{3.} $f_X=x+\tfrac12$; $\E(X)=\tfrac13+\tfrac14=\tfrac7{12}$; $\E(X^2)=\tfrac14+\tfrac16=\tfrac5{12}$; $\Var(X)=\tfrac5{12}-\tfrac{49}{144}=\tfrac{11}{144}$.
$\E(XY)=\int_0^1\!\int_0^1(x^2y+xy^2)\,dx\,dy=\tfrac16+\tfrac16=\tfrac13$.
$\Cov=\tfrac13-\tfrac{49}{144}=-\tfrac1{144}$; $\rho=-\tfrac1{11}$.}
\end{turno}

\begin{ojo}
\mal{Usar $\E(XY)=\E(X)\E(Y)$ para calcular la covariancia.}
\bien{Eso s\'olo vale si son independientes, y entonces la covariancia es cero. En general $\E(XY)$ se calcula con la conjunta.}
\mal{Dividir entre las variancias: $\rho=\Cov/(\Var(X)\Var(Y))$.}
\bien{Se divide entre las \textbf{desviaciones est\'andar}: $\rho=\Cov/(\sigma_X\sigma_Y)$. Con variancias se obtiene $\rho^2$ si arriba va la covariancia al cuadrado.}
\mal{Aceptar un $\rho$ fuera de $[-1,1]$ o una variancia negativa.}
\bien{Son se\~nales seguras de un error de cuentas. Revisa $\E(X^2)$ y los cuadrados.}
\end{ojo}
